\uuid{SGuu}
\titre{Pgcd}
\auteur{bodin ; croizat ; sacre}
\organisation{exo7}
\licence{CC BY-NC-SA 4.0}
\datecreate{2022-03-08}
\niveau{L1}
\module{Algèbre}
\chapitre{Arithmétique dans Z}
\sousChapitre{Pgcd, ppcm, algorithme d'Euclide}
\difficulte{2}
\theme{pgcd}
\exo7qcmid{hilisit.arithmetique.005}

\contenu{
\begin{question}
\qclassification{103.03}
    Soit $a=bq+r$ la division euclidienne de $a$ par $b$.
    Quelle condition définit le reste $r$ ?
    \begin{answers}
        \bad{$0 \le r < a$}
        \good{$0 \le r < b$}
        \bad{$0 \le r \le q$}
        \bad{$0 \le r < q$}
    \end{answers}
    \begin{explanations}
    Dans la division euclidienne, on a $0 \le r \le b-1$, c’est-à-dire $0 \le r < b$ puisque $r$ est un entier.
    Cela permet d'avoir l'unicité du quotient $q$ et du reste $r$.
    \end{explanations}
\end{question}
}
